\sheet{1}{The ECU Sees Only Signals}{Chapter~1 (The ECU Sees Only Signals)}

\begin{goals}
\begin{itemize}
  \item Classify the engine sensor signals (continuous/discrete time,
    analog/binary amplitude, deterministic/stochastic, periodic/aperiodic)
    and know their characteristic time scales.
  \item \emph{Bonus (optional):} describe sensor signals quantitatively:
    mean, RMS, extreme values, autocorrelation as a numerical test for
    periodicity.
  \item Extract a first piece of information from a raw signal: crank
    angle and engine speed from the 60-2 tooth timing, including the
    effect of the timer resolution.
  \item Synchronise crank and cam signal to determine the $720^\circ$
    phase of the engine cycle (cylinder identification).
\end{itemize}
\end{goals}

\section*{Background: the 60-2 trigger wheel}

The crank wheel of the virtual engine has 60 tooth positions ($6^\circ$ each),
two of which are missing. The Hall sensor delivers a binary signal; the
rising edge of tooth $k=0,\dots,57$ occurs at $\theta=6k^\circ$ (modulo
$360^\circ$), the reference mark is the first rising edge after the gap.

\begin{center}
\begin{tikzpicture}[x=0.16cm,y=0.7cm,font=\footnotesize]
  % teeth 54..57, gap, teeth 0..3 drawn from -36 deg to +24 deg
  \draw[->,gray] (-40,0) -- (30,0) node[right,black] {$\theta$};
  \foreach \a in {-36,-30,-24,-18,0,6,12,18}
    \draw[thick,dspblue] (\a,0) -- (\a,1) -- (\a+3,1) -- (\a+3,0);
  \draw[thick,dspblue] (-40,0) -- (-36,0) (-33,0) -- (-30,0) (-27,0) -- (-24,0)
       (-21,0) -- (-18,0) (-15,0) -- (0,0) (3,0)--(6,0) (9,0)--(12,0) (15,0)--(18,0) (21,0)--(28,0);
  \node[below] at (-18,0) {$342^\circ$};
  \node[below] at (0,0) {$0^\circ/360^\circ$};
  \node[below] at (12,0) {$12^\circ$};
  \node[above] at (-18,1) {57};
  \node[above] at (-24,1) {56};
  \node[above] at (0,1) {0};
  \node[above] at (6,1) {1};
  \draw[<->] (-24,1.6) -- node[above] {$T_{56}$ ($6^\circ$)} (-18,1.6);
  \draw[<->] (-18,1.6) -- node[above] {$T_{\text{gap}}$ ($18^\circ$)} (0,1.6);
  \node[align=center] at (-7.5,0.5) {gap};
\end{tikzpicture}
\end{center}

The tooth period $T_k$ is measured from rising edge to rising edge. Because one
tooth corresponds to $6^\circ = 1/60$ revolution, the engine speed in rpm is
\begin{equation}
  n = \frac{60\,\text{s/min}}{60\,T_k} = \frac{1}{T_k/\text{s}}\ \text{rpm},
  \qquad
  n = \frac{3}{T_{\text{gap}}/\text{s}}\ \text{rpm (gap)},
  \qquad
  n = \frac{60}{T_{\text{rev}}/\text{s}}\ \text{rpm (revolution)}.
  \label{eq:rpm}
\end{equation}


\section*{Background: crankshaft, camshaft and the $720^\circ$ cycle}

A four-stroke engine needs \emph{two} crankshaft revolutions for one working
cycle (intake, compression, power, exhaust), but every valve opens only
once per cycle. The camshaft is therefore driven by a timing belt (or chain)
with a ratio of $2:1$: its sprocket has twice as many teeth as the crankshaft
sprocket, so it turns once while the crankshaft turns twice. The crank angle
$\theta$ is counted over the full cycle, $0^\circ\le\theta<720^\circ$, which
is exactly one camshaft revolution ($360^\circ$ cam angle). The 60-2 wheel on
the crankshaft repeats after $360^\circ$ crank angle, the single segment of
the cam phase wheel only after $720^\circ$.

\begin{center}
\begin{tikzpicture}[font=\footnotesize,>={Stealth[length=4pt]}]
  % ---------------- (a) mechanism, front view (schematic) ----------------
  \begin{scope}
    % 60-2 trigger wheel on the crankshaft (58 teeth, 3 deg each)
    \draw[dspblue] (0,0) circle[radius=1.0];
    \foreach \i in {0,...,57} {
      \pgfmathsetmacro\a{150-6*\i}
      \fill[dspblue] (\a:1.0) -- (\a:1.12)
        arc[start angle=\a, end angle=\a-3, radius=1.12] -- ({\a-3}:1.0) -- cycle;
    }
    % cylinder, piston, connecting rod, crank throw
    \draw[thick] (-0.48,1.45) -- (-0.48,4.0) (0.48,1.45) -- (0.48,4.0);
    \filldraw[fill=gray!25] (-0.62,4.0) rectangle (0.62,4.35);
    \filldraw[fill=gray!40] (-0.42,2.49) rectangle (0.42,3.09);
    \draw[very thick,dspgray] (0,0) -- (0.45,0.536);
    \draw[very thick,dspgray] (0.45,0.536) -- (0,2.689);
    \fill (0.45,0.536) circle[radius=0.05] (0,2.689) circle[radius=0.05];
    % sprockets and timing belt (ratio 2:1)
    \filldraw[fill=white,thick] (0,0) circle[radius=0.55];
    \fill (0,0) circle[radius=0.06];
    \filldraw[fill=white,thick] (0,5.6) circle[radius=1.1];
    \draw[line width=1.6pt,white]
      (0.547,-0.054) arc[start angle=-5.64, end angle=-174.36, radius=0.55]
      -- (-1.095,5.492) arc[start angle=185.64, end angle=-5.64, radius=1.1]
      -- cycle;
    \draw[line width=0.8pt]
      (0.547,-0.054) arc[start angle=-5.64, end angle=-174.36, radius=0.55]
      -- (-1.095,5.492) arc[start angle=185.64, end angle=-5.64, radius=1.1]
      -- cycle;
    % cam phase wheel segment (90 deg cam = 180 deg crank)
    \fill[dspblue] (20:1.2) ++(0,5.6) arc[start angle=20, end angle=110, radius=1.2]
      -- ($(0,5.6)+(110:1.32)$) arc[start angle=110, end angle=20, radius=1.32] -- cycle;
    % cam lobe, valve
    \filldraw[fill=gray!40] (0.3,5.6) arc[start angle=0, end angle=180, radius=0.3]
      to[out=-90,in=180] (0,5.08) to[out=0,in=-90] cycle;
    \fill (0,5.6) circle[radius=0.06];
    \draw[thick] (0,5.08) -- (0,4.05);
    \fill (-0.22,3.95) -- (0.22,3.95) -- (0.05,4.05) -- (-0.05,4.05) -- cycle;
    % Hall sensors
    \filldraw[fill=dspgray] (225:1.2) ++(-0.28,-0.12) rectangle ++(0.22,0.24);
    \filldraw[fill=dspgray] ($(0,5.6)+(60:1.4)$) rectangle ++(0.24,0.22);
    % labels
    \draw[gray] (-1.1,-0.85) -- (-1.6,-1.2) node[below,black,align=center]
      {crank sensor};
    \draw[gray] ($(0,5.6)+(62:1.6)$) -- (1.4,7.4) node[above,black] {cam sensor};
    \node[align=right,anchor=east] at (-1.25,0.3)
      {crankshaft,\\ $z$ teeth,\\ 60-2 wheel};
    \node[align=right,anchor=east] at (-1.25,5.6)
      {camshaft,\\ $2z$ teeth,\\ phase segm.};
    \node[anchor=west,align=left] at (1.1,2.6) {timing\\ belt $2:1$};
    % rotation arrows
    \draw[->,thick,dspblue] (60:1.4) arc[start angle=60, end angle=10, radius=1.4];
    \node[dspblue,anchor=west,align=left] at (1.45,0.65)
      {2 rev.\\ $=720^\circ$};
    \draw[->,thick,dspblue] ($(0,5.6)+(-20:1.5)$) arc[start angle=-20, end angle=-60, radius=1.5];
    \node[dspblue,anchor=west,align=left] at (1.55,4.35)
      {1 rev.\\ $=360^\circ$};
    \node at (0,-1.9) {(a)};
  \end{scope}
  % ---------------- (b) one working cycle over the crank angle -------------
  \begin{scope}[shift={(4.85,0.4)},x=0.0108cm,y=1cm]
    % row labels
    \foreach \y/\t in {5.4/crankshaft,4.6/camshaft,3.8/cylinder 1,2.9/crank signal,2.1/cam signal}
      \node[anchor=east] at (-8,\y) {\t};
    % crankshaft revolutions
    \draw (0,5.15) rectangle node {1st revolution} (360,5.65);
    \draw (360,5.15) rectangle node {2nd revolution} (720,5.65);
    % camshaft revolution
    \draw (0,4.35) rectangle node {1 revolution ($360^\circ$ cam angle)} (720,4.85);
    % strokes of cylinder 1 (combustion TDC at 120 deg)
    \draw[fill=gray!15] (0,3.55) rectangle node {compr.} (120,4.05);
    \draw[fill=orange!35] (120,3.55) rectangle node {power} (300,4.05);
    \draw[fill=gray!15] (300,3.55) rectangle node {exhaust} (480,4.05);
    \draw[fill=dspblue!15] (480,3.55) rectangle node {intake} (660,4.05);
    \draw[fill=gray!15] (660,3.55) rectangle (720,4.05);
    % crank signal: 60-2 teeth, repeats every 360 deg
    \foreach \r in {0,360} \foreach \k in {0,...,57}
      \fill[dspblue] (\r+6*\k,2.65) rectangle (\r+6*\k+3,3.15);
    \draw[dspblue] (0,2.65) -- (720,2.65);
    % cam signal: HIGH for 0..180 deg, repeats every 720 deg
    \draw[thick,dspblue] (0,1.85) -- (0,2.35) -- (180,2.35) -- (180,1.85) -- (720,1.85);
    % axis
    \draw[->] (0,1.5) -- (760,1.5) node[right] {$\theta$};
    \foreach \x in {0,120,180,360,480,720}
      \draw (\x,1.45) -- (\x,1.55) node[pos=0,below] {$\x^\circ$};
    % TDCs of cylinder 1
    \draw[dashed,red!70!black] (120,1.5) -- (120,4.25);
    \draw[dashed,red!70!black] (480,1.5) -- (480,4.25);
    \node[red!70!black,align=center] at (120,0.6) {TDC\\ ignition};
    \node[red!70!black,align=center] at (480,0.6) {TDC\\ gas exchange};
    \node at (360,-0.25) {(b)};
  \end{scope}
\end{tikzpicture}
\end{center}
\noindent (a)~Schematic front view of one cylinder with crank drive and
overhead camshaft; the timing belt turns the camshaft once per two crankshaft
revolutions. (b)~One working cycle of cylinder~1 over the crank angle: the
crank wheel pattern and the piston position repeat after $360^\circ$, the
strokes and the cam signal only after $720^\circ$.
%-------------------------------------------------------------------------
\begin{exercise}{Signal classification and time scales}{\Paper}
\begin{tasks}
  \item Classify the five engine signals of Table~1.2 of the script (crank,
    cam, knock, ion current, lambda) with respect to
    (i) time: continuous/discrete, (ii) amplitude: analog/binary (digital),
    (iii) deterministic/stochastic, (iv) periodic/quasi-periodic/aperiodic.
    Justify every entry in one line. Note that the answer may differ for the
    physical sensor output and for the representation inside the ECU.
  \item Discuss the borderline cases: (i) Is the crank signal periodic
    during an acceleration? Is it periodic in another independent
    variable? (ii) Why does the script call the knock signal
    \emph{quasi-periodic}, although a single knock event is a transient?
    Introduce the notion of a signal whose \emph{statistics} are periodic
    in the crank angle.
  \item Fill in the following table for an engine speed of \SI{3000}{rpm}:
    crank revolution period, tooth period and tooth frequency, duration of the
    gap (rising to rising edge), cam period, cam HIGH time, interval between
    two combustions (firing frequency), duration of $1^\circ$ crank angle,
    time window of a knock onset ($5^\circ$--$25^\circ$ after TDC), number of
    \SI{6.5}{kHz} knock oscillations within one decay time
    $\tau=\SI{0.8}{ms}$.
  \item Crank and cam: the crankshaft turns $720^\circ$ per engine cycle,
    the camshaft $360^\circ$. Why is the crank signal alone not sufficient to
    decide whether the piston of cylinder~1 is in its compression or in its
    exhaust stroke? After switching on the ignition with the engine
    cranking at \SI{200}{rpm}, what is the worst-case time until the ECU knows
    (i) the crank angle modulo $360^\circ$, (ii) the full $720^\circ$ phase,
    if the cam level is evaluated at $\theta = 60^\circ$ after the reference
    mark?
\end{tasks}
\end{exercise}

\begin{solution}
a) Classification (physical signal $\to$ inside the ECU):

\begin{center}
\small
\begin{tabularx}{\textwidth}{@{}l>{\raggedright\arraybackslash}p{2.6cm}>{\raggedright\arraybackslash}p{2.3cm}>{\raggedright\arraybackslash}p{3.3cm}>{\raggedright\arraybackslash}X@{}}
\toprule
Signal & time & amplitude & predictability & periodicity \\
\midrule
crank & CT $\to$ event times $t_k$ & binary (Hall) & deterministic up to small
  speed fluctuations & periodic in $\theta$ ($360^\circ$); in $t$ only at constant speed \\
cam & CT $\to$ level sampled at tooth events & binary & deterministic &
  periodic in $\theta$ ($720^\circ$) \\
knock & CT $\to$ DT after ADC & analog $\to$ quantised & stochastic (noise,
  random knock onset, amplitude, phase) & quasi-periodic: deterministic parts
  repeat every $180^\circ$, knock is random \\
ion current & CT $\to$ DT & analog & partly deterministic (shape of the
  peaks), cycle-to-cycle variation & aperiodic per event; event pattern
  repeats per $180^\circ$ \\
lambda & CT $\to$ DT (slow, 10--100\,Hz) & analog, but practically
  two-valued (0.1/0.9\,V) & slowly varying, disturbed & aperiodic; in the
  closed loop a limit cycle (quasi-periodic, $\approx$\SI{1}{Hz}) \\
\bottomrule
\end{tabularx}
\end{center}
Crank and cam are physically continuous-time voltages with two levels; the
ECU, however, does not sample them but time-stamps their edges with a
capture timer -- the information is in the \emph{sequence of edge times}
$t_k$, i.e.\ a discrete-time signal indexed by the tooth number.
All analog signals become discrete-time and quantised (= digital) after the ADC.

b) (i) During an acceleration the tooth period changes from tooth to tooth,
so $x(t)=x(t+T)$ holds for no $T$: the crank signal is not periodic in time.
It is, however, \emph{exactly} periodic in the crank angle,
$x_c(\theta)=x_c(\theta+360^\circ)$, independently of the speed -- the reason
why many engine algorithms work in the angle domain (Sheet~4).
(ii) Each knock burst is a transient (damped oscillation), its onset,
amplitude and phase are random and it does not occur in every cycle. But the
\emph{probability} of a burst, the valve impacts, the combustion vibration
and the noise level repeat every $180^\circ$ (each TDC): the mean value and the
autocorrelation of the signal are periodic functions of the crank angle.
Such a signal is called \emph{cyclostationary}. ``Quasi-periodic'' in the
script means exactly this.

c) At $n=\SI{3000}{rpm}$ ($\SI{50}{rev/s}$, $\omega=\SI{18000}{\degree/s}$):
\begin{center}
\small
\begin{tabular}{@{}lll@{}}
\toprule
quantity & formula & value \\
\midrule
crank revolution & $60/n$ & \SI{20}{ms} (\SI{50}{Hz}) \\
tooth period / frequency & $T_{\text{rev}}/60$ & \SI{333}{\micro s} / \SI{3}{kHz} \\
gap (rising to rising) & $3\,T_z$ & \SI{1}{ms} \\
engine cycle = cam period & $2\,T_{\text{rev}}$ & \SI{40}{ms} (\SI{25}{Hz}) \\
cam HIGH time & $180^\circ$ & \SI{10}{ms} \\
combustion interval & $720^\circ/4 = 180^\circ$ & \SI{10}{ms} (firing frequency \SI{100}{Hz} = 2nd order) \\
$1^\circ$ crank angle & $1/\omega$ & \SI{55.6}{\micro s} \\
knock onset window & $5^\circ\dots25^\circ$ ATDC & \SI{0.28}{ms} \dots \SI{1.39}{ms} after TDC \\
knock oscillations in $\tau$ & $f_1\tau$ & $6.5\cdot 0.8 = 5.2$ \\
\bottomrule
\end{tabular}
\end{center}
The signals thus span four decades of time scales: knock (\SI{100}{\micro s}),
teeth (\SI{300}{\micro s}), combustions (\SI{10}{ms}), lambda (\SI{1}{s}).

d) The crank wheel looks identical in both revolutions of a cycle: at
$\theta=120^\circ$ (compression TDC of cylinder~1, ignition) and at
$\theta=480^\circ$ (gas-exchange TDC of cylinder~1, overlap) the crank signal
is the same. Only the camshaft, which turns once per cycle, distinguishes the
two revolutions. At \SI{200}{rpm} one revolution lasts \SI{300}{ms}.
(i) Worst case: the ECU starts just after the reference mark and must wait for
the next gap \emph{and} needs one normal tooth period before it to apply the
ratio test: almost one revolution, $\approx\SI{300}{ms}$.
(ii) Then another $60^\circ$ to the cam evaluation:
$420^\circ \approx \SI{350}{ms}$. Before that, the ECU can only use
``wasted spark'' (fire both cylinders 1 and 4 at every crank TDC) and must not
inject sequentially.
\end{solution}

%-------------------------------------------------------------------------
\begin{exercise}{60-2 crank wheel and engine speed}{\Wokwi}
Open the Wokwi project \codefile{wokwi/sheet01\_crank/} (Arduino Uno, engine
simulator chip, logic analyzer, two LEDs, push button). \texttt{eng:CRANK}
is connected to pin~2 (external interrupt INT0), \texttt{eng:CAM} to pin~4.
The ISR \texttt{crank\_isr()} is called at every rising crank edge and
time-stamps it. With \texttt{TIMEBASE\_TIMER1 = 1} the time stamp is the
16-bit Timer1 counter (prescaler 8, \SI{0.5}{\micro s} per tick), with 0 it
is \texttt{micros()} (resolution \SI{4}{\micro s} on the Uno).
\begin{tasks}
  \item \emph{Preparation (paper).} Derive \eqref{eq:rpm}. A timer of resolution
    $\Delta$ quantises the tooth period. Show that the resulting quantisation
    step of the per-tooth speed is $\Delta n \approx n^2\Delta$ (with $n$ in rpm
    and $\Delta$ in s) and evaluate it for $\Delta=\SI{4}{\micro s}$ and
    \SI{0.5}{\micro s} at 600, 3000 and \SI{7000}{rpm}. Compare with the
    amplitude of the 1\,\% speed ripple of the model. The gap is detected by
    the ratio test $T_k > 2\,T_{k-1}$: estimate the angular acceleration
    (in rpm/s) at which the test would fail (missed gap, or false gap), and
    compare with realistic engine accelerations
    ($\lesssim\SI{10000}{rpm/s}$).
  \item Complete the ISR: gap detection, tooth counter, per-revolution
    statistics, TDC LED (pin 12) on for $30^\circ$ starting at each TDC
    edge. Which tooth numbers belong to $\theta=120^\circ$ and $300^\circ$?
    Check with the logic analyzer (D0 = crank, D2 = TDC output of the chip,
    D3 = your LED) and report the column \texttt{teeth} and
    \texttt{sync\_err}.
  \item Complete the speed computation in \texttt{loop()} (per revolution,
    last tooth, min/max per tooth). At \SI{3000}{rpm}, compare
    \texttt{rpm\_rev} with the spread \texttt{rpm\_min}\dots\texttt{rpm\_max}
    for both time bases. Press the button to dump 120 tooth periods, copy
    only this block and plot the per-tooth speed (column \texttt{rpm}), e.g.\
    in a spreadsheet or with your CSV plot program of Bonus Exercise~0.B1\,d)
    in JSLinux. Which part of the
    ``jitter'' is information, which part is measurement error?
  \item Move the speed slider (600\dots\SI{7000}{rpm}) and repeat c) at
    \SI{7000}{rpm}. Why is printing one CSV line per \emph{tooth} not an
    option (estimate the UART load)?
\end{tasks}
\end{exercise}

\begin{solution}
a) One tooth $=6^\circ=\frac{1}{60}$\,rev, so the speed is
$\frac{1}{60\,T_k}$\,rev/s $=\frac{60}{60\,T_k}$\,rpm $=1/T_k$ (in rpm with $T_k$
in s). The gap covers $18^\circ$, a revolution $360^\circ$, hence $3/T_{\text{gap}}$
and $60/T_{\text{rev}}$. With $n=1/T$: $\left|\frac{\dd n}{\dd T}\right|=\frac{1}{T^2}=n^2$,
so a period error of one tick changes $n$ by $\Delta n\approx n^2\Delta$:

\begin{center}
\small
\begin{tabular}{@{}rrrrr@{}}
\toprule
$n$ / rpm & $T_k$ & $\Delta n$ ($\Delta=\SI{4}{\micro s}$) & $\Delta n$ ($\Delta=\SI{0.5}{\micro s}$) & ripple ($\pm1\,\%$) \\
\midrule
600  & \SI{1667}{\micro s} & \SI{1.4}{rpm} & \SI{0.2}{rpm} & $\pm$\SI{6}{rpm} \\
3000 & \SI{333}{\micro s}  & \SI{36}{rpm}  & \SI{4.5}{rpm} & $\pm$\SI{30}{rpm} \\
7000 & \SI{143}{\micro s}  & \SI{196}{rpm} & \SI{24.5}{rpm}& $\pm$\SI{70}{rpm} \\
\bottomrule
\end{tabular}
\end{center}
The relative resolution $\Delta n/n=\Delta/T_k$ degrades linearly with speed. With
\texttt{micros()} the ripple is buried in the quantisation above
$\approx\SI{2500}{rpm}$; with the \SI{0.5}{\micro s} timer it is resolved over the
whole range. Real ECUs use capture timers with $\le\SI{100}{ns}$ resolution.

Gap test: let $\omega_1$ be the speed during the tooth before the gap,
$\omega_2$ the mean speed across the gap. Then $T_{\text{gap}}/T_{\text{prev}}
=3\omega_1/\omega_2<2$ requires $\omega_2>1.5\,\omega_1$ within about two
tooth periods ($\approx12^\circ/\omega_1$):
$\dot\omega\approx 0.5\,\omega_1^2/12^\circ$, i.e.\ $\dot n\approx n^2/4$ (rpm/s,
$n$ in rpm) -- \SI{90000}{rpm/s} at \SI{600}{rpm}. A false gap
($T_k/T_{k-1}>2$ for a normal tooth) needs the speed to halve within one tooth,
$\dot n\approx -n^2/2$, \SI{-180000}{rpm/s} at \SI{600}{rpm}. Both are an order
of magnitude beyond real engines, so threshold 2 (the geometric mean of 1 and 3
would be 1.73) is robust in normal operation. Critical is only the
start (cranking at \SI{200}{rpm} with $\pm$30\,\% speed fluctuations from the
compression strokes, limit $\dot n\approx\SI{10000}{rpm/s}$), where ECUs add
plausibility checks (tooth count 58 between gaps, cam edge position).

b) Key part of the solution
(\codefile{solutions/wokwi/sheet01\_crank/sketch.ino}):
\inputcode[firstline=52,lastline=99]{solutions/wokwi/sheet01_crank/sketch.ino}
$120^\circ$ and $300^\circ$ correspond to the rising edges of teeth 20 and 50.
Every revolution reports \texttt{teeth} $=58$ edges (57 normal teeth + the gap
edge) and \texttt{sync\_err} $=0$. In the logic analyzer the LED rises
a few microseconds (interrupt latency) after the rising edge of tooth 20 (50),
i.e.\ practically together with the \SI{50}{\micro s} TDC pulse of the chip, which
is emitted exactly at $120^\circ$/$300^\circ$. The LED marks the
crank-angle \emph{event}; it is lit for 5 teeth = $30^\circ$.

c) Expected values at \SI{3000}{rpm} (obtained with an offline model of exactly this
measurement, \codefile{solutions/wokwi/sheet01\_crank/tooth\_sim.c}, over 200
revolutions; in Wokwi rare additional outliers of a few \si{\micro s} can occur
when the crank interrupt has to wait for the Timer0 or UART interrupt):
\begin{center}
\small
\begin{tabular}{@{}lccc@{}}
\toprule
time base & rpm\_rev & per-tooth min\dots max & per-tooth std \\
\midrule
\texttt{micros()}, \SI{4}{\micro s} & 2999.4\dots3000.0 & 2941\dots3049 & \SI{25.7}{rpm} \\
Timer1, \SI{0.5}{\micro s} & 2999.85 & 2967\dots3030 & \SI{21.7}{rpm} \\
\bottomrule
\end{tabular}
\end{center}
The per-revolution speed is very precise (one tick on \SI{20}{ms}:
$\Delta n/n=0.5/40000$) and slightly \emph{below} 3000, because the mean of
$1/(1+\varepsilon\sin2\theta)$ is $1+\varepsilon^2/2$
($n\varepsilon^2/2=\SI{0.15}{rpm}$). The per-tooth values with Timer1 show the
\emph{real} speed fluctuation: $\pm1\,\%$, two periods per revolution, std
$30/\sqrt2=\SI{21.2}{rpm}$ -- this is information (in a real engine the
torque pulses of the individual combustions; used e.g.\ for misfire
detection). With \texttt{micros()} the values jump in steps of $\approx\SI{36}{rpm}$
(332, 336\,\si{\micro s}) -- a measurement error of the same size as the
information. The plotted dump shows a clean sine with 2 periods per
58 teeth for Timer1 and a staircase-like, noisy version for \texttt{micros()}.

d) At \SI{7000}{rpm} the model gives per-tooth std \SI{94}{rpm} (micros, range
6757\dots7143) vs.\ \SI{51}{rpm} (Timer1, 6920\dots7092, ripple
$\pm\SI{70}{rpm}$); the per-revolution value stays within
$\pm\SI{3}{rpm}$ resp.\ \SI{0.5}{rpm}. When the slider is moved, the speed
changes abruptly in the simulator; the gap is still found (ratio 3 vs.\ 1,
see a)) as long as the jump does not fall exactly into the gap.
UART: at \SI{7000}{rpm} there are $7000\cdot58/60\approx6800$ teeth/s; a line
such as \texttt{1234,57,286,6993.0} has $\approx20$ characters, i.e.\
\SI{136000}{chars/s} -- twelve times the capacity of 115200\,baud
(\SI{11520}{chars/s}). Hence statistics per revolution and block dumps from RAM.
\end{solution}

%-------------------------------------------------------------------------
\begin{exercise}{Cam/crank synchronisation}{\Challenge}
Extend your sketch such that it determines the $720^\circ$ phase from the cam
signal (\texttt{eng:CAM}, HIGH for $0^\circ\le\theta<180^\circ$).
\begin{tasks}
  \item Why should the cam level not be read at the reference edge
    ($\theta=0^\circ$), and which tooth is a good choice instead?
  \item Implement \texttt{phase} ($0^\circ$ or $360^\circ$), toggle it at every
    gap, and check it at every revolution (\texttt{cam\_err}). Light the LED on
    pin~11 only at the compression TDC of cylinder~1 and verify it with the
    logic analyzer (D1 = cam, D4 = LED).
  \item Give the crank angle and the number of the cylinder that fires at each of
    the four TDC events.
\end{tasks}
\end{exercise}

\begin{solution}
a) The cam edge at $\theta=0^\circ$ coincides with the reference edge of the
crank signal (in the simulator both are written in the same event, on a real
engine the cam edge has a tolerance of several degrees and drifts with cam
phasing). Reading the level there is a race. The level must be read in the
middle of a cam state, e.g.\ at tooth 10 ($60^\circ$ or $420^\circ$): HIGH $\Rightarrow$
first revolution (phase $0^\circ$), LOW $\Rightarrow$ second revolution (phase
$360^\circ$). Any tooth with $\ge 10^\circ$ distance to $0^\circ$ and $180^\circ$ works.

b) See the listing in Exercise~1.2 (line 21 and lines 43--48): at every gap
\texttt{phase = 360 - phase}, at tooth 10 the cam level is compared with the
predicted phase; \texttt{cam\_err} stays 0 and the printed \texttt{phase}
alternates 0, 360, 0, \dots{} The green LED is on only every second time
the red TDC LED lights up, i.e.\ once per engine cycle (every \SI{40}{ms} at
\SI{3000}{rpm}).

c) $\theta=\text{phase}+6^\circ k$; firing order 1-3-4-2:
\begin{center}
\small
\begin{tabular}{@{}cccc@{}}
\toprule
phase & tooth & crank angle & cylinder \\
\midrule
$0^\circ$   & 20 & $120^\circ$ & 1 \\
$0^\circ$   & 50 & $300^\circ$ & 3 \\
$360^\circ$ & 20 & $480^\circ$ & 4 \\
$360^\circ$ & 50 & $660^\circ$ & 2 \\
\bottomrule
\end{tabular}
\end{center}
\end{solution}

\begin{deliverables}
1.1: completed classification table and time-scale table. 1.2: sketch,
screenshot of the logic analyzer (gap, TDC, LED), table of rpm statistics for
both time bases at 3000 and \SI{7000}{rpm}, plot of a tooth-period dump. 1.3
(optional): extended sketch, logic analyzer screenshot with cam and
cylinder-1 LED. Bonus 1.B1 (optional): program, table of statistics and
periods, short interpretation.
\end{deliverables}

%=========================================================================
\section*{Bonus: JSLinux (optional)}
\begingroup
\setcounter{exercise}{0}\renewcommand{\theexercise}{B\arabic{exercise}}
For interested students: the following exercise analyses all engine signals
offline in JSLinux (set up as in Bonus Exercise~0.B1).

\begin{exercise}{Signal zoo}{\JSLinux}
The program \codefile{bonus/jslinux/sheet01/signal\_zoo.c} generates \SI{160}{ms}
(four engine cycles) of the crank, cam, knock and ion signals at
\SI{3000}{rpm}, $f_s=\SI{25}{kHz}$, and \SI{20}{s} of a lambda sensor signal
($f_s=\SI{200}{Hz}$) for a slowly oscillating $\lambda(t)$. Additionally, a crank
signal during a speed ramp \SIrange{2000}{4000}{rpm} is generated.
\begin{tasks}
  \item Complete \texttt{stats()}: mean, RMS value, minimum and maximum. Run
    the program and explain the values of the crank and cam signals
    analytically. Why is the mean of the knock signal not zero, although the
    knock sensor is AC coupled?
  \item The normalised autocorrelation
    \[ r[m] = \frac{\hat R[m]}{\hat R[0]},\qquad
       \hat R[m]=\frac{1}{N-m}\sum_{n=0}^{N-1-m}(x[n]-\bar x)(x[n+m]-\bar x) \]
    is close to 1 at lags $m$ where the signal resembles a shifted copy of
    itself. Complete \texttt{autocorr()} and \texttt{period\_lag()} (first local
    maximum after the first zero crossing that reaches at least 90\,\% of the
    largest value). Report the estimated period and $r$ for all six signals.
  \item Interpret the results: Which signals are periodic, which only
    quasi-periodic? Why does the crank signal show a strong peak already at
    half a revolution? What happens to the periodicity of the crank signal
    during the ramp -- and which lag is still detected? Why would the
    $1/N$-normalised (biased) estimate be misleading here?
\end{tasks}
\end{exercise}

\begin{solution}
a) Output of the reference solution \codefile{solutions/bonus/jslinux/sheet01/signal\_zoo.c}:
\begin{lstlisting}[style=shell,numbers=none]
signal              mean       rms       min       max
crank             0.4808    0.6934    0.0000    1.0000
cam               0.2500    0.5000    0.0000    1.0000
knock             0.0433    0.1234   -1.1454    1.5144
ion               0.0993    0.2399   -0.1269    1.0440
lambda            0.5004    0.6195    0.1000    0.9000
\end{lstlisting}
For a 0/1 signal with duty cycle $d$: mean $=d$, RMS $=\sqrt{d}$.
Crank: $d=58\cdot3^\circ/360^\circ=0.4833$ (measured 0.4808 because a HIGH
phase of $3^\circ=\SI{167}{\micro s}$ is only 4.17 samples long, sampling
rounds), RMS $=\sqrt{0.4808}=0.6934$. Cam: $d=180/720=0.25$, RMS $=0.5$.
The knock signal contains the low-frequency combustion vibration, a positive
half sine of amplitude \SI{0.2}{V} over $60^\circ$ after every TDC:
$0.2\cdot\frac{2}{\pi}\cdot\frac{60}{180}=\SI{0.042}{V}$ -- this is the measured
mean. (A real piezo sensor is AC-coupled; the model represents the signal
after a high-pass with a cut-off far below the combustion frequency.)
The large extremes ($-1.15/+1.51$\,V) stem from knock events.

b), c) Result of the period search:
\begin{lstlisting}[style=shell,numbers=none]
crank 3000   first peak: lag  250 =   10.00 ms r = 0.929 | max: lag  500 r = 0.988
crank ramp   first peak: lag    8 =    0.32 ms r = 0.348 | max: lag    8 r = 0.348
cam          first peak: lag 1000 =   40.00 ms r = 1.000 | max: lag 1000 r = 1.000
knock        first peak: lag  244 =    9.76 ms r = 0.418 | max: lag  500 r = 0.438
ion          first peak: lag  250 =   10.00 ms r = 0.992 | max: lag  250 r = 0.992
lambda       first peak: lag  251 = 1255.00 ms r = 0.981 | max: lag  500 r = 0.982
\end{lstlisting}
\begin{itemize}
  \item \emph{Cam}: exactly periodic, $r=1.000$ at \SI{40}{ms} = one engine cycle.
  \item \emph{Crank}: the true period is one revolution (\SI{20}{ms}, $r=0.988$;
    not exactly 1 because of the speed ripple and because the period is
    500.03 samples, not an integer). A shift by $180^\circ$ maps 58 teeth onto
    teeth except near the gap, therefore $r(250)=0.929$ is almost as high:
    the gap, which carries the position information, is only a small part of
    the signal energy. With the 90\,\% rule the detector reports
    $180^\circ$ -- a reminder that a period detector must be designed for the
    feature of interest. (Small peaks at multiples of 8.33 samples come from
    the teeth.)
  \item \emph{Crank during the ramp}: the tooth period shrinks from 12.5 to
    6.25 samples; no long-lag peak survives; only the tooth-scale peak at lag
    8 ($r=0.35$) remains. In time the signal is not periodic -- in the
    angle domain it would be (Sheet~4).
  \item \emph{Knock}: $r\approx0.42$ at $\approx\SI{10}{ms}$
    ($180^\circ$): only the deterministic parts (combustion vibration, valve
    impacts at $8$\,kHz) repeat; noise and the random knock bursts
    (random onset and phase) do not correlate. Quasi-periodic, cyclostationary.
  \item \emph{Ion current}: $r=0.992$ at $180^\circ$ -- in the model all
    cylinders produce the same deterministic ion curve; on a real engine
    cycle-to-cycle variations make each event an individual (aperiodic)
    transient.
  \item \emph{Lambda}: period \SI{1.25}{s} = $1/\SI{0.8}{Hz}$, $r=0.98$; the
    switching sensor turns the small $\lambda$ oscillation into an almost
    rectangular signal.
\end{itemize}
The biased estimator ($1/N$ instead of $1/(N-m)$) multiplies $r[m]$ by
$(N-m)/N$; at lag 1000 of a 4000-sample record this is a factor $0.75$, so the
cam signal would show $r=0.75$ although it is perfectly periodic, and shorter
lags would always be favoured. Key part of the solution:
\inputcode[firstline=48,lastline=80]{solutions/bonus/jslinux/sheet01/signal_zoo.c}
\end{solution}
\endgroup
